\chapter{Non-Local Gate Synthesis and Multi-Controlled Operations}
\label{ch:nonlocal_gate_synthesis}

\begin{quote}
\textit{``When the control qubit sits in dilution refrigerator Alpha and the target qubit in dilution refrigerator Beta, direct unitary evolution is physically impossible. The compiler must synthesize non-local interaction through quantum teleportation and phase kickback.''}
\end{quote}

Once the partitioning engine has divided the quantum circuit across multiple physical QPUs, all gates that cross partition boundaries must be transformed into physically realizable operations. Because direct transverse coupling cannot span macroscopic inter-QPU distances, non-local quantum interactions are synthesized through distributed sub-circuits that consume pre-shared entanglement (EPR pairs), local entangling gates, and classical communication.

This chapter presents the mathematical theory, algebraic derivations, and circuit architectures for non-local gate synthesis. We analyze the canonical two-qubit non-local gates (CNOT and Controlled-Phase), detail the Cat-entangler/disentangler mechanism, provide complete multi-QPU decompositions for three-qubit Toffoli (CCX) gates across arbitrary partitions, and formalize the general $n$-qubit Multi-Controlled-X (MCX) distributed synthesis algorithms implemented in the DQC compiler.

% ==============================================================================
\section{Canonical Non-Local Two-Qubit Gates}
\label{sec:canonical_2q_synthesis}
% ==============================================================================

The foundational building block of distributed quantum logic is the non-local Controlled-NOT (CNOT) operation between a control qubit $\ket{\psi_c}$ on QPU $A$ and a target qubit $\ket{\psi_t}$ on QPU $B$.

\subsection{Algebraic Derivation of the Eisert Tele-Gate Protocol}
Let the initial quantum state before gate execution be:
\begin{equation}
    \ket{\Psi_0} = \ket{\psi_c}_A \otimes \ket{\psi_t}_B = (\alpha \ket{0}_A + \beta \ket{1}_A) \otimes (\gamma \ket{0}_B + \delta \ket{1}_B).
    \label{eq:initial_bipartite_state}
\end{equation}

To execute a remote CNOT without moving the primary qubits, QPUs $A$ and $B$ share an ancillary Bell pair:
\begin{equation}
    \ket{\Phi^+}_{e_A e_B} = \frac{1}{\sqrt{2}} \left( \ket{0}_{e_A} \ket{0}_{e_B} + \ket{1}_{e_A} \ket{1}_{e_B} \right).
\end{equation}
The four-qubit system state is $\ket{\Psi_1} = \ket{\Psi_0} \otimes \ket{\Phi^+}_{e_A e_B}$.

\begin{figure}[htbp]
    \centering
    \begin{tikzpicture}[scale=0.92, >=stealth]
        % QPU Alpha Box
        \draw[rounded corners=3mm, fill=blue!5, draw=darkblue, line width=1pt] (-5.5, -2.8) rectangle (-0.5, 2.2);
        \node[font=\bfseries\color{darkblue}] at (-3, 1.8) {QPU Alpha (Sender)};
        
        % QPU Beta Box
        \draw[rounded corners=3mm, fill=blue!5, draw=darkblue, line width=1pt] (0.5, -2.8) rectangle (5.5, 2.2);
        \node[font=\bfseries\color{darkblue}] at (3, 1.8) {QPU Beta (Receiver)};
        
        % Horizontal Lines
        \draw[line width=1pt] (-5, 1.0) node[left] {$q_c$} -- (-1, 1.0) -- (0, 1.0) -- (5, 1.0) node[right] {$q_c$};
        \draw[line width=1pt] (-5, -0.2) node[left] {$e_A$} -- (-1, -0.2);
        \draw[line width=1pt] (1, -0.2) -- (5, -0.2) node[right] {$e_B$};
        \draw[line width=1pt] (1, -1.8) node[left] {$q_t$} -- (5, -1.8) node[right] {$q_t$};
        
        % Shared EPR link
        \draw[<->, line width=1.5pt, color=quantumviolet, dashed] (-1, -0.2) to [bend left=25] 
            node[above, font=\scriptsize\bfseries] {Shared $\ket{\Phi^+}$} (1, -0.2);
            
        % Step 1: Local CNOT on Alpha (qc -> eA)
        \node[draw, circle, fill=black, inner sep=2pt] (c1) at (-4, 1.0) {};
        \node[draw, circle, fill=white, inner sep=3.5pt] (t1) at (-4, -0.2) {};
        \draw[line width=1pt] (t1.north) -- (t1.south);
        \draw[line width=1pt] (t1.east) -- (t1.west);
        \draw[line width=1pt] (c1) -- (t1);
        \node[font=\tiny, left] at (-4.1, 0.4) {1. CX};
        
        % Step 2: Measure eA on Alpha in Z basis
        \draw[fill=white, draw=black, line width=0.8pt] (-2.5, -0.5) rectangle (-1.7, 0.1);
        \node[font=\tiny] at (-2.1, -0.2) {$\mathcal{M}_Z$};
        \node[font=\tiny, above] at (-2.1, 0.1) {$m_1$};
        
        % Step 3: Local CNOT on Beta (eB -> qt)
        \node[draw, circle, fill=black, inner sep=2pt] (c2) at (2, -0.2) {};
        \node[draw, circle, fill=white, inner sep=3.5pt] (t2) at (2, -1.8) {};
        \draw[line width=1pt] (t2.north) -- (t2.south);
        \draw[line width=1pt] (t2.east) -- (t2.west);
        \draw[line width=1pt] (c2) -- (t2);
        \node[font=\tiny, right] at (2.1, -1.0) {3. CX};
        
        % Step 4: Hadamard on eB
        \draw[fill=white, draw=black, line width=0.8pt] (2.8, -0.5) rectangle (3.4, 0.1);
        \node[font=\tiny] at (3.1, -0.2) {$H$};
        
        % Step 5: Measure eB in X basis (H + Mz)
        \draw[fill=white, draw=black, line width=0.8pt] (3.8, -0.5) rectangle (4.4, 0.1);
        \node[font=\tiny] at (4.1, -0.2) {$\mathcal{M}_Z$};
        \node[font=\tiny, above] at (4.1, 0.1) {$m_2$};
        
        % Classical signaling lines
        % m1 to target fixup (X correction)
        \draw[double, ->, line width=0.8pt, color=darkblue] (-2.1, -0.5) -- (-2.1, -2.3) -- (4.2, -2.3) -- (4.2, -1.8);
        \node[draw, fill=white, font=\tiny, inner sep=1.5pt] at (4.2, -1.8) {$X^{m_1}$};
        
        % m2 back to control fixup (Z correction)
        \draw[double, ->, line width=0.8pt, color=compilerteal] (4.1, 0.1) -- (4.1, 1.5) -- (-0.8, 1.5) -- (-0.8, 1.0);
        \node[draw, fill=white, font=\tiny, inner sep=1.5pt] at (-0.8, 1.0) {$Z^{m_2}$};
    \end{tikzpicture}
    \caption{Circuit architecture of the Eisert-Jacobs-Papadopoulos-Plenio (EJPP) non-local CNOT gate. The gate transfers the control excitation to QPU $B$ via entangling ancilla $e_A$, applies local CNOT on QPU $B$, and kicks back the phase correlation to $q_c$ via classical bit $m_2$.}
    \label{fig:ejpp_circuit_diagram}
\end{figure}

The execution proceeds through five synchronized steps:
\begin{enumerate}
    \item \textbf{Local CNOT on QPU $A$ ($q_c \to e_A$):}
    \begin{equation}
        \ket{\Psi_2} = \frac{1}{\sqrt{2}} \left[ \alpha \ket{0}_{q_c} \ket{0}_{e_A} \ket{0}_{e_B} + \beta \ket{1}_{q_c} \ket{1}_{e_A} \ket{0}_{e_B} + \alpha \ket{0}_{q_c} \ket{1}_{e_A} \ket{1}_{e_B} + \beta \ket{1}_{q_c} \ket{0}_{e_A} \ket{1}_{e_B} \right] \otimes \ket{\psi_t}_B.
    \end{equation}
    
    \item \textbf{Computational Measurement of $e_A$ on QPU $A$:}
    Measuring $e_A$ yields classical bit $m_1 \in \{0, 1\}$ and projects the remaining state into:
    \begin{equation}
        \ket{\Psi_3} = \alpha \ket{0}_{q_c} \ket{m_1}_{e_B} \ket{\psi_t}_B + \beta \ket{1}_{q_c} \ket{m_1 \oplus 1}_{e_B} \ket{\psi_t}_B.
    \end{equation}
    Notice that the state of $e_B$ on QPU $B$ is now entangled with the control qubit $q_c$ on QPU $A$.
    
    \item \textbf{Local CNOT on QPU $B$ ($e_B \to q_t$):}
    Applying local CNOT treats $e_B$ as the effective control qubit targeting $q_t$:
    \begin{equation}
        \ket{\Psi_4} = \alpha \ket{0}_{q_c} \ket{m_1}_{e_B} (\sigma_x^{m_1} \ket{\psi_t}_B) + \beta \ket{1}_{q_c} \ket{m_1 \oplus 1}_{e_B} (\sigma_x^{m_1 \oplus 1} \ket{\psi_t}_B).
    \end{equation}
    
    \item \textbf{Hadamard and Measurement of $e_B$ on QPU $B$ (Phase Kickback):}
    To decouple ancilla $e_B$ from the data register, QPU $B$ applies a Hadamard gate $H$ to $e_B$ and measures in the computational basis, yielding bit $m_2 \in \{0, 1\}$. Expanding the state in the Hadamard basis reveals:
    \begin{equation}
        \ket{\Psi_5} = \frac{1}{\sqrt{2}} \left[ \alpha \ket{0}_{q_c} (\sigma_x^{m_1}\ket{\psi_t}) + (-1)^{m_2} \beta \ket{1}_{q_c} (\sigma_x^{m_1 \oplus 1}\ket{\psi_t}) \right].
    \end{equation}
    
    \item \textbf{Classical Feedforward Fix-Up:}
    \begin{itemize}
        \item QPU $A$ sends bit $m_1$ to QPU $B$. QPU $B$ applies local bit-flip $\sigma_x^{m_1}$ to $q_t$.
        \item QPU $B$ sends bit $m_2$ to QPU $A$. QPU $A$ applies local phase-flip $\sigma_z^{m_2}$ to $q_c$.
    \end{itemize}
\end{enumerate}

After these corrections, the final joint state is:
\begin{equation}
    \ket{\Psi_{\text{final}}} = \alpha \ket{0}_{q_c} \ket{\psi_t}_B + \beta \ket{1}_{q_c} (\sigma_x \ket{\psi_t}_B) \equiv \text{CNOT}(q_c, q_t) \left( \ket{\psi_c} \otimes \ket{\psi_t} \right).
\end{equation}
The operation is mathematically deterministic, completing a universal non-local CNOT using **exactly 1 EPR pair, 2 local CNOTs, and 2 classical bits** (1 forward, 1 backward).

% ==============================================================================
\section{The Cat-Entangler and Cat-Disentangler Framework}
\label{sec:cat_entangler}
% ==============================================================================

When a single control qubit on QPU $P_0$ must simultaneously control multiple target qubits distributed across $k$ distinct QPUs $\{P_1, P_2, \dots, P_k\}$ (a common pattern in multi-controlled gates, QFT, and syndrome extraction), executing independent point-to-point tele-gates requires $k$ separate EPR pairs and induces deep serial dependencies.

The DQC compiler optimizes this pattern using the **Cat-Entangler / Cat-Disentangler** paradigm. Instead of independent pairwise teleportations, the compiler synthesizes a distributed Greenberger-Horne-Zeilinger (GHZ) ``Cat'' state spanning all target QPUs.

\begin{figure}[htbp]
    \centering
    \begin{tikzpicture}[scale=0.9, >=stealth]
        % Horizontal wire labels
        \node[font=\footnotesize] at (-4.5, 3) {QPU 0 ($q_c$)};
        \node[font=\footnotesize] at (-4.5, 1.8) {QPU 0 ($a_0$)};
        \node[font=\footnotesize] at (-4.5, 0.6) {QPU 1 ($a_1$)};
        \node[font=\footnotesize] at (-4.5, -0.6) {QPU 2 ($a_2$)};
        \node[font=\footnotesize] at (-4.5, -1.8) {QPU 3 ($a_3$)};
        
        % Horizontal wires
        \draw[line width=1pt] (-3, 3) -- (7, 3);
        \draw[line width=1pt] (-3, 1.8) -- (7, 1.8);
        \draw[line width=1pt] (-3, 0.6) -- (7, 0.6);
        \draw[line width=1pt] (-3, -0.6) -- (7, -0.6);
        \draw[line width=1pt] (-3, -1.8) -- (7, -1.8);
        
        % Cat Entangler Stage
        \draw[fill=blue!10, draw=darkblue, rounded corners=2mm, line width=1pt] (-2.5, -2.1) rectangle (-0.2, 2.1);
        \node[font=\footnotesize\bfseries\color{darkblue}, rotate=90] at (-1.35, 0) {Cat-Entangler};
        
        % Coupling from qc to Cat state
        \node[draw, circle, fill=black, inner sep=2pt] at (0.5, 3) {};
        \node[draw, circle, fill=white, inner sep=3.5pt] (t0) at (0.5, 1.8) {};
        \draw[line width=1pt] (t0.north) -- (t0.south);
        \draw[line width=1pt] (t0.east) -- (t0.west);
        \draw[line width=1pt] (0.5, 3) -- (t0);
        
        % Parallel Target Operations
        \draw[fill=compilerteal!15, draw=compilerteal, rounded corners=2mm, line width=1pt] (1.5, -2.1) rectangle (3.5, 1.0);
        \node[font=\small\bfseries\color{compilerteal}, align=center] at (2.5, -0.55) {Parallel Local\\Target Gates\\($U_{t_1}, U_{t_2}, U_{t_3}$)};
        
        % Cat Disentangler Stage
        \draw[fill=purple!10, draw=quantumviolet, rounded corners=2mm, line width=1pt] (4.2, -2.1) rectangle (6.5, 2.1);
        \node[font=\footnotesize\bfseries\color{quantumviolet}, rotate=90] at (5.35, 0) {Cat-Disentangler};
        
        % Annotations
        \draw[decorate, decoration={brace, amplitude=5pt}, line width=0.8pt] (-2.5, 2.3) -- (-0.2, 2.3);
        \node[above=6pt, font=\tiny] at (-1.35, 2.3) {1. Distribute GHZ};
        
        \draw[decorate, decoration={brace, amplitude=5pt}, line width=0.8pt] (1.5, 1.2) -- (3.5, 1.2);
        \node[above=6pt, font=\tiny] at (2.5, 1.2) {2. Execute in Parallel};
        
        \draw[decorate, decoration={brace, amplitude=5pt}, line width=0.8pt] (4.2, 2.3) -- (6.5, 2.3);
        \node[above=6pt, font=\tiny] at (5.35, 2.3) {3. Disentangle \& Correct};
    \end{tikzpicture}
    \caption{The Cat-Entangler / Cat-Disentangler operational architecture. A single control qubit $q_c$ entangles into a distributed Cat state $(\ket{00\dots0} + \ket{11\dots1})$, which simultaneously controls gates on QPUs 1, 2, and 3 in parallel $\mathcal{O}(1)$ depth before being disentangled via reverse LOCC.}
    \label{fig:cat_entangler_architecture}
\end{figure}

The Cat-state distribution requires $k$ EPR pairs to create an $(k+1)$-qubit shared state:
\begin{equation}
    \ket{\text{Cat}}_{a_0 a_1 \dots a_k} = \frac{1}{\sqrt{2}} \left( \ket{0}^{\otimes (k+1)} + \ket{1}^{\otimes (k+1)} \right).
\end{equation}
Once established, non-local control operations across all $k$ target QPUs execute simultaneously in parallel circuit depth $\mathcal{O}(1)$, reducing circuit latency from $\mathcal{O}(k)$ to a single broadcast round.

% ==============================================================================
\section{Distributed Multi-Controlled Gates (CCX / Toffoli)}
\label{sec:distributed_toffoli}
% ==============================================================================

The three-qubit Toffoli (CCX) gate is ubiquitous in reversible arithmetic (Draper adders, Grover diffusion operators, Shor's modular exponentiation). When its three qubits are distributed across multiple QPUs, synthesis complexity depends entirely on the spatial partition configuration:

\subsection{Partition Topology 1: Control-Target Colocated (2 QPUs)}
Assume control $c_0$ resides on QPU $A$, while control $c_1$ and target $t$ reside on QPU $B$. 
In this configuration, target QPU $B$ can perform a local Toffoli if control $c_0$ can be temporarily teleported or shared.
Using the Margolus (relative-phase) Toffoli decomposition, the circuit requires **exactly 1 non-local CNOT** (1 EPR pair) if relative phase accumulation can be corrected locally, or **2 non-local CNOTs** (2 EPR pairs) for an exact Toffoli.

\subsection{Partition Topology 2: Completely Disjoint (3 QPUs)}
The most challenging configuration occurs when control $c_0 \in \text{QPU } A$, control $c_1 \in \text{QPU } B$, and target $t \in \text{QPU } C$. No two qubits share a physical chip.

\begin{theorembox}{Minimum Entanglement Bound for 3-QPU Toffoli}
Any deterministic implementation of a three-qubit Toffoli gate $\text{CCX}(c_0, c_1, t)$ partitioned across three completely disjoint quantum processors ($c_0 \in P_A, c_1 \in P_B, t \in P_C$) requires a minimum of:
\begin{equation}
    N_{\text{EPR}} \ge 3 \quad \text{EPR pairs},
    \label{eq:toffoli_epr_lower_bound}
\end{equation}
and at least 4 rounds of classical communication.
\end{theorembox}

\begin{figure}[htbp]
    \centering
    \begin{tikzpicture}[scale=0.9, >=stealth]
        % Three QPU horizontal bands
        \draw[rounded corners=2mm, fill=blue!5, draw=darkblue, line width=0.8pt] (-5, 2.2) rectangle (5, 3.8);
        \node[font=\bfseries\color{darkblue}, left] at (-5.2, 3) {QPU A};
        \draw[line width=1pt] (-4.5, 3) node[left] {$c_0$} -- (4.5, 3) node[right] {$c_0$};
        
        \draw[rounded corners=2mm, fill=green!5, draw=compilerteal, line width=0.8pt] (-5, 0.2) rectangle (5, 1.8);
        \node[font=\bfseries\color{compilerteal}, left] at (-5.2, 1) {QPU B};
        \draw[line width=1pt] (-4.5, 1) node[left] {$c_1$} -- (4.5, 1) node[right] {$c_1$};
        
        \draw[rounded corners=2mm, fill=purple!5, draw=quantumviolet, line width=0.8pt] (-5, -1.8) rectangle (5, -0.2);
        \node[font=\bfseries\color{quantumviolet}, left] at (-5.2, -1) {QPU C};
        \draw[line width=1pt] (-4.5, -1) node[left] {$t$} -- (4.5, -1) node[right] {$t$};
        
        % Distributed T and T^\dagger gates
        % Step 1: Remote CNOT between A and B
        \draw[<->, line width=1.5pt, color=darkblue] (-3, 3) -- node[midway, fill=white, font=\tiny] {Telegate 1} (-3, 1);
        \draw[fill=white, draw=black] (-2, 0.7) rectangle (-1.4, 1.3);
        \node[font=\tiny] at (-1.7, 1) {$T^\dagger$};
        
        % Step 2: Remote CNOT between B and C
        \draw[<->, line width=1.5pt, color=compilerteal] (-0.5, 1) -- node[midway, fill=white, font=\tiny] {Telegate 2} (-0.5, -1);
        \draw[fill=white, draw=black] (0.5, -1.3) rectangle (1.1, -0.7);
        \node[font=\tiny] at (0.8, -1) {$T$};
        
        % Step 3: Remote CNOT between A and C
        \draw[<->, line width=1.5pt, color=quantumviolet] (2.2, 3) -- node[midway, fill=white, font=\tiny] {Telegate 3} (2.2, -1);
        \draw[fill=white, draw=black] (3.2, -1.3) rectangle (3.8, -0.7);
        \node[font=\tiny] at (3.5, -1) {$T^\dagger$};
    \end{tikzpicture}
    \caption{Distributed synthesis of a Toffoli gate across three physically isolated QPUs ($P_A, P_B, P_C$). Non-local telegates mediate intermediate phase accumulation, consuming 3 EPR pairs across the tripartite network.}
    \label{fig:tripartite_toffoli}
\end{figure}

The DQC compiler synthesizes this tripartite Toffoli by decomposing the unitary into the Clifford+$T$ basis:
\begin{align}
    \text{CCX} &= (H \otimes \mathbb{I} \otimes \mathbb{I}) \cdot \text{CNOT}(c_1, t) \cdot T_t^\dagger \cdot \text{CNOT}(c_0, t) \cdot T_t \nonumber \\
    &\quad \cdot \text{CNOT}(c_1, t) \cdot T_t^\dagger \cdot \text{CNOT}(c_0, t) \cdot (T_{c_1} \otimes T_t \otimes H) \dots
\end{align}
By applying telegate commutation passes, the compiler merges redundant tele-CNOT operations, strictly achieving the theoretical lower bound of 3 EPR pairs.

% ==============================================================================
\section{General $n$-Qubit Multi-Controlled-X (MCX) Synthesis}
\label{sec:mcx_synthesis}
% ==============================================================================

For general multi-controlled gates $\text{C}^k\text{NOT}$ with $k \ge 3$ control qubits distributed across $M$ QPUs, direct Clifford+$T$ decomposition without ancillae scales exponentially as $\mathcal{O}(2^k)$ T-gates. 

To achieve polynomial efficiency, the DQC compiler implements **Distributed Ancilla Borrowing**. Every QPU contains idle computation qubits that are temporarily unentangled during the execution window. The compiler borrows these clean idle qubits as dirty ancillae, constructing a distributed V-ladder network:

\begin{equation}
    \text{EPR}_{\text{MCX}}(k, M) \le 2(M - 1) + 4 \cdot \text{Cut}(k),
    \label{eq:mcx_cost_formula}
\end{equation}
where $\text{Cut}(k)$ is the number of control wires severed by the partition boundary.

This synthesis algorithm guarantees that any arbitrary quantum program can be distributed over a multi-QPU network with bounded communication overhead, establishing the mathematical transformation rules executed by Pass 3 of our MLIR compilation pipeline.
